% 分类目录：ml
\documentclass[12pt, a4paper]{article}
\usepackage[margin=2.5cm]{geometry}
\usepackage{ctex}
\usepackage{amsmath, amssymb}
\usepackage{listings}
\usepackage{xcolor}
\usepackage{tabularx}
\usepackage{hyperref}

\lstset{
    language=Python,
    basicstyle=\ttfamily\small,
    keywordstyle=\color{blue},
    commentstyle=\color{green!50!black},
    stringstyle=\color{red},
    showstringspaces=false,
    breaklines=true,
    numbers=left,
    numberstyle=\tiny\color{gray},
    frame=single
}

\title{知识点：因子分析 (Factor Analysis)}
\author{}
\date{}

\begin{document}
\maketitle

\section{核心定义}
\textbf{中文定义}：因子分析（FA）是一种降维和数据分析技术，用于研究观测变量之间的相关性，并通过少数几个不可观测的隐变量（潜因子）来解释这种相关性。它假设每个观测变量都是由若干个“公共因子”和唯一的“特殊因子”线性组合而成的。

\textbf{英文定义}：Factor Analysis (FA) is a dimensionality reduction and data analysis technique used to model the correlations between observed variables using a smaller number of unobserved latent variables called factors. It assumes that each observed variable is a linear combination of several "common factors" and a unique "specific factor."

\section{数学模型}
假设观测变量向量为 $\mathbf{X} = (x_1, x_2, \dots, x_p)^T$，模型表示为：
$$ \mathbf{X} = \mathbf{\Lambda} \mathbf{F} + \mathbf{\epsilon} $$
其中：
\begin{itemize}
    \item $\mathbf{F} = (f_1, f_2, \dots, f_m)^T$ 是 $m$ 个公共因子 ($m < p$)。
    \item $\mathbf{\Lambda}$ 是 $p \times m$ 的 \textbf{因子载荷矩阵} (Factor Loading Matrix)，$\lambda_{ij}$ 表示第 $i$ 个变量在第 $j$ 个因子上的载荷（相关程度）。
    \item $\mathbf{\epsilon}$ 是特殊因子，包含了变量不能被公共因子解释的部分。
\end{itemize}

\section{核心概念：方差分解}
对于观测变量 $x_i$，其方差可以分解为两部分：
\begin{enumerate}
    \item \textbf{共同性 (Communality)}：由公共因子解释的方差部分，即载荷平方和 $h_i^2 = \sum_{j=1}^m \lambda_{ij}^2$。
    \item \textbf{特殊方差 (Specific Variance)}：公共因子无法解释的部分。
\end{enumerate}

\section{因子旋转 (Factor Rotation)}
获取初始载荷矩阵后，通常载荷分布比较分散，难以解释。通过“旋转”载荷阵，可以使每个变量在少数因子上有大载荷，在其他因子上载荷接近 0（简单结构原则）。
\begin{itemize}
    \item \textbf{正交旋转 (Varimax)}：最常用，保持因子之间互不相关，使因子载荷的方差最大化。
    \item \textbf{斜交旋转 (Promax)}：允许因子之间存在相关性，往往能得到更符合现实逻辑的解释。
\end{itemize}

\section{因子分析 vs 主成分分析 (PCA)}
\begin{table}[h]
\centering
\begin{tabularx}{\textwidth}{|l|X|X|}
\hline
\textbf{维度} & \textbf{主成分分析 (PCA)} & \textbf{因子分析 (FA)} \\
\hline
\textbf{目标} & 寻找原始变量的线性组合以解释最大方差。 & 寻找潜在的结构（潜因子）以解释相关性。 \\
\hline
\textbf{本质} & $Y = AX$，由变量推导出成分。 & $X = \Lambda F + \epsilon$，由因子推导出变量。 \\
\hline
\textbf{方差处理} & 包含所有总方差。 & 仅关注公共方差（共同性）。 \\
\hline
\textbf{解释性} & 结果客观，但物理意义不明确。 & 结果主观，但具有明确的潜在物理或商业含义。 \\
\hline
\end{tabularx}
\end{table}

\section{主要步骤：因子性检验}
在进行分析前，需要确认数据是否适合做因子分析：
\begin{itemize}
    \item \textbf{KMO 检验}：取值越接近 1 越好，通常 > 0.6 表示适合。
    \item \textbf{Bartlett 球形检验}：检验相关阵是否为单位阵，若 $p < 0.05$ 则拒绝原假设，表示变量间存在相关性。
\end{itemize}

\section{关键代码片段 (Scikit-learn)}
\begin{lstlisting}
from sklearn.decomposition import FactorAnalysis
import pandas as pd

# 1. 初始化模型，设置提取因子数
fa = FactorAnalysis(n_components=3, rotation='varimax')

# 2. 拟合数据并转换
fa_result = fa.fit_transform(X)

# 3. 查看载荷矩阵
loadings = pd.DataFrame(fa.components_.T, columns=['F1', 'F2', 'F3'])
print(loadings)
\end{lstlisting}

\section{深度面试题}

\textbf{问：为什么说因子分析的结果比 PCA 更具主观色彩？}\\
答：因为因子分析允许对载荷阵进行旋转，且因子的命名和解释很大程度上依赖于研究者的业务背景。此外，研究者需要预先设定提取的因子个数，这会直接影响最终的共同性解释。

\textbf{问：共同性 (Communality) 低于 0.4 意味着什么？}\\
答：这意味着该观测变量不能很好地被现有的公共因子所解释，该变量可能带有较多噪音，或者它反映的是一个独立于其他变量的、未被捕获的新维度。

\section{参考文献}
\begin{enumerate}
    \item Child, D. (2006). \textit{The Essentials of Factor Analysis}. Bloomsbury Academic.
    \item Costello, A. B., \& Osborne, J. (2005). \textit{Best Practices in Exploratory Factor Analysis}. Practical Assessment, Research \& Evaluation.
    \item Harman, H. H. (1976). \textit{Modern Factor Analysis}. University of Chicago Press.
\end{enumerate}

\end{document}
